% ===== PRÉAMBULE CANONIQUE — à coller VERBATIM en tête de CHAQUE .tex =====
\documentclass[11pt,a4paper]{article}
\usepackage[utf8]{inputenc}\usepackage[T1]{fontenc}\usepackage{lmodern}
\usepackage[french]{babel}
\usepackage{amsmath,amssymb,mathtools}\usepackage{icomma}
\usepackage{siunitx}\sisetup{locale=FR}  % \qty{}{}, \si{}, \ang{}
\usepackage{xcolor}\usepackage{colortbl}
\usepackage{tikz}\usepackage{pgfplots}\pgfplotsset{compat=1.18}
\usepackage[siunitx,europeanresistors]{circuitikz} % circuits (élec.)
\usepackage[most]{tcolorbox}\usepackage{enumitem}\usepackage{array,booktabs}
\usepackage[a4paper,margin=2.2cm]{geometry}
\usetikzlibrary{arrows.meta,calc,angles,quotes,positioning,patterns,%
  backgrounds,shapes.geometric,decorations.pathmorphing,decorations.markings,intersections}
\definecolor{primary}{RGB}{222,49,99}\definecolor{primarydark}{RGB}{150,22,55}
\definecolor{defcol}{RGB}{0,114,178}\definecolor{methcol}{RGB}{52,92,148}
\definecolor{excol}{RGB}{0,135,90}\definecolor{attncol}{RGB}{200,40,55}
\tcbset{boxbase/.style={enhanced,breakable,boxrule=0.9pt,arc=2.5pt,
  left=3mm,right=3mm,top=2.3mm,bottom=2.3mm,fonttitle=\bfseries,coltitle=white}}
% --- boîtes TUTORIEL (noms EXACTS, ne pas renommer) ---
\newtcolorbox{cle}[1][]{boxbase,colback=defcol!6,colframe=defcol,
  title={\textsc{À retenir}\ifx\relax#1\relax\relax\else\textmd{ : \itshape #1}\fi}}
\newtcolorbox{methode}[1][]{boxbase,colback=methcol!7,colframe=methcol,sharp corners,
  title={\textsc{Méthode}\ifx\relax#1\relax\relax\else\textmd{ --- \itshape #1}\fi}}
\newtcolorbox{exemplebox}[1][]{boxbase,colback=excol!5,colframe=excol,
  title={\textsc{Exemple}\ifx\relax#1\relax\relax\else\textmd{ : \itshape #1}\fi}}
\newtcolorbox{piege}[1][]{boxbase,colback=attncol!6,colframe=attncol,
  title={\textsc{Piège}\ifx\relax#1\relax\relax\else\textmd{ : \itshape #1}\fi}}
\newtcolorbox{remarque}[1][]{boxbase,colback=black!4,colframe=black!45,
  title={\textsc{Remarque}\ifx\relax#1\relax\relax\else\textmd{ : \itshape #1}\fi}}
% --- boîtes EXERCICE / CORRIGÉ (noms EXACTS, auto-comptées) ---
\newtcolorbox[auto counter]{exo}[1][]{boxbase,colback=primary!4,colframe=primary,
  title={\textsc{Exercice}~\thetcbcounter\ifx\relax#1\relax\relax\else\textmd{ --- \itshape #1}\fi}}
\newtcolorbox[auto counter]{corrige}[1][]{boxbase,colback=excol!4,colframe=excol,
  title={\textsc{Corrigé}~\thetcbcounter\ifx\relax#1\relax\relax\else\textmd{ --- \itshape #1}\fi}}
\newcommand{\vv}[1]{\vec{#1}}\newcommand{\ds}{\displaystyle}
\newcommand{\R}{\mathbb{R}}\newcommand{\N}{\mathbb{N}}\newcommand{\Z}{\mathbb{Z}}
% (un \usepackage{fancyhdr}+en-tête est possible ; il est ignoré côté web)
% =========================================================================
\begin{document}

\begin{center}
{\large\bfseries\color{excol} Thème 3 — Corrigé Facile}\\[-2pt]
{\color{excol}\rule{5cm}{1pt}}
\end{center}

\begin{corrige}[compter sur une onde stationnaire]
On compte directement sur le schéma (les points marquent les nœuds).
\begin{enumerate}[label=\alph*)]
  \item \textbf{3 fuseaux} (trois « ventres » successifs).
  \item \textbf{3 ventres} : autant que de fuseaux.
  \item \textbf{4 nœuds} : un mode à $n$ fuseaux possède $n+1$ nœuds ($3+1=4$),
        les deux extrémités comprises.
\end{enumerate}
\end{corrige}

\begin{corrige}[distances sur une onde stationnaire]
On retient : nœud--nœud $=\lambda/2$, nœud--ventre voisin $=\lambda/4$.
\begin{enumerate}[label=\alph*)]
  \item $\dfrac{\lambda}{2}=\dfrac{0{,}6}{2}=\SI{0.3}{\metre}.$
  \item $\dfrac{\lambda}{4}=\dfrac{0{,}6}{4}=\SI{0.15}{\metre}.$
\end{enumerate}
\end{corrige}

\begin{corrige}[longueurs d'onde des modes (deux bouts fixes)]
On applique $\lambda_n=2L/n$ avec $L=\SI{1.5}{\metre}$, donc $2L=\SI{3}{\metre}$.
\begin{enumerate}[label=\alph*)]
  \item $\lambda_1=\dfrac{2L}{1}=\SI{3}{\metre}.$
  \item $\lambda_2=\dfrac{2L}{2}=L=\SI{1.5}{\metre}.$
  \item $\lambda_3=\dfrac{2L}{3}=\dfrac{3}{3}=\SI{1}{\metre}.$
\end{enumerate}
\end{corrige}

\begin{corrige}[les harmoniques sont des multiples]
Les fréquences propres sont des multiples entiers du fondamental : $f_n=n\,f_1$.
\begin{enumerate}[label=\alph*)]
  \item $f_2=2\times 40=\SI{80}{\hertz}.$
  \item $f_3=3\times 40=\SI{120}{\hertz}.$
  \item $f_5=5\times 40=\SI{200}{\hertz}.$
\end{enumerate}
\end{corrige}

\begin{corrige}[corde à une extrémité libre]
Le bout libre est un ventre : le mode fondamental correspond à un quart de longueur d'onde,
d'où $\lambda_1=4L$.
\begin{enumerate}[label=\alph*)]
  \item $\lambda_1=4L=4\times 0{,}6=\SI{2.4}{\metre}.$
  \item $f_1=\dfrac{c}{\lambda_1}=\dfrac{240}{2{,}4}=\SI{100}{\hertz}.$
\end{enumerate}
\end{corrige}

\end{document}
